% One-page letter / statement template. Matches the CV letterhead.
\documentclass[11pt, letterpaper]{article}
\usepackage[T1]{fontenc}
\usepackage[utf8]{inputenc}
\usepackage{charter}
\usepackage[top=0.55in, bottom=0.45in, left=0.8in, right=0.8in]{geometry}
\usepackage[dvipsnames]{xcolor}
\definecolor{linkblue}{HTML}{1155CC}
\usepackage[colorlinks=true, urlcolor=linkblue, linkcolor=linkblue]{hyperref}
\usepackage{parskip}
\pagestyle{empty}
\setlength{\parindent}{0pt}
\linespread{1.0}

\newcommand{\letterhead}{%
  \begin{center}
    {\large\bfseries Rodrigo França Chaves}\\[2pt]
    {\small
      Chicago, IL \enspace$\cdot$\enspace
      \href{mailto:rchaves@uchicago.edu}{rchaves@uchicago.edu} \enspace$\cdot$\enspace
      +1 312 217 5526 \enspace$\cdot$\enspace
      \href{https://rodrigofch7.github.io}{rodrigofch7.github.io}
    }
  \end{center}
  \vspace{2pt}
  {\rule{\linewidth}{0.8pt}}
  \vspace{10pt}
}

\begin{document}
\letterhead
\begin{center}{\bfseries\large Statement of Interest: Professor Shaoda Wang}\\[2pt]
{\small Algorithmic regulation and environmental enforcement targeting}\end{center}
\vspace{6pt}

The premise of this project is the one I find most interesting in applied regulation: enforcement capacity is finite, so the targeting rule is itself a policy instrument, and a rule based on measured pollution alone ignores that identical emissions cause very different damage depending on location, transport, and who lives downwind. I would like to work on it because I have built the analytical machinery this question needs, and because I have learned, in a different setting, how easily that machinery misleads when the targeting criterion is chosen carelessly.

The closest precedent in my own work is a horizontal-equity audit of New York City's property tax roll. Enforcement there faces the same structural problem: an assessment agency cannot review 1.1 million properties, so it needs a defensible rule for deciding which ones merit attention. Comparing each property to a citywide benchmark is useless, because a Manhattan cooperative and a Queens two-family fall into the same column. I built a peer-group construction that measures each property only against genuinely comparable ones, defined by borough, building class, tax class, decade built, size, and value band, with a five-level fallback so rare building types still receive a comparison set. Roughly 43 percent sit more than 15 percent from their peer-group median. The essential caveat is that both the peer definition and the threshold were choices we made, so that figure measures dispersion within comparable groups rather than proven misassessment. I designed the labeling methodology, built the data pipeline, and trained the models, using 138 leakage-free features and gradient boosting. Tightening the value bands from quintiles to ventiles cut the within-group price spread from \$475,000 to \$48,000, which cost a little accuracy and bought labels that meant something. That trade is, I think, the central methodological issue in algorithmic enforcement: the label definition is the model.

I also bring the spatial side. At the World Bank I worked on a 100-meter gridded panel of Dakar built from Google Open Buildings and GHSL rasters, estimating construction responses to a new transit corridor with difference-in-differences. The most consequential thing I found was a measurement problem rather than a result: the building-detection threshold had been set near twice the value its source paper recommends for the region, discarding real structures most heavily in dense and informal areas. Since that bias varies by neighborhood, it does not difference out. For a project estimating damages that depend on population density and exposure, I would expect similar hazards in the spatial layer, and I would look for them early.

On method and identification, my first-authored article in \textit{Utilities Policy} uses staggered difference-in-differences with propensity-matched controls to estimate the health effects of transferring Brazilian water systems to private operators. I wrote the paper, assembled the data, and wrote the code, and I am comfortable reporting mixed results when the data give them.

I work fluently in Python, R, and Stata, with geospatial workflows in QGIS and geopandas and version-controlled, reproducible pipelines. What I would contribute is the construction and validation of the analysis data, the feature and targeting work, and a persistent skepticism about whether a rule that performs well on our own labels is measuring social harm or merely reproducing how the existing system already behaves.

\end{document}
